% ----------------------------------------------------------------------
%  3.1.1. Классификация конденсированных систем
% ----------------------------------------------------------------------
\begin{orientir}
Почему атомы держатся вместе? Ответ дают законы квантовой механики,
которые определяют строение электронных оболочек атомов. В этом разделе
мы: вспомним, чем конденсированные тела отличаются от газов (3.1.1);
разберём строение атомов и периодическую систему (3.1.2); поймём,
какие силы связывают атомы и как оценить энергию связи (3.1.3--3.1.4);
подробно рассмотрим четыре типа связи~— ковалентную, металлическую,
ионную и ван-дер-ваальсову~— и то, какие кристаллические структуры
каждая из них порождает (3.1.5).
\end{orientir}

\subsection{Классификация конденсированных систем}\label{ssec:condensed}

\term{Конденсированные системы}~— твёрдые тела и жидкости (вдали от
критической точки). От газов они отличаются принципиально.
\begin{itemize}
\item \emph{Частицы касаются друг друга}: расстояния между ними равны их
диаметрам, тогда как в газе при атмосферном давлении~— примерно в десять
раз больше. Поэтому сжимаемость конденсированных тел в миллион раз меньше,
чем у газов.
\item \emph{Тепловое движение~— колебания}, а не поступательное движение.
Лишь изредка частица перескакивает в соседний пустой узел или междоузлие
(глава~2). Колебания квантуются, и при низких температурах это приводит к
интересным квантовым явлениям (п.~\ref{sec:thermal}).
\item \emph{Собственный объём}: конденсированное тело удерживается
внутренними силами сцепления без внешних сил, а газ занимает весь
предоставленный объём.
\end{itemize}

Только при сверхвысоких давлениях (тысячи и миллионы атмосфер),
сравнимых с внутренним давлением сил отталкивания, удаётся сжать твёрдое
тело на 20--30\,\% по объёму. При этом вещество сначала переходит в
металлическое, а затем в плазменное состояние, и твёрдое тело как
таковое перестаёт существовать.

По возрастанию степени упорядоченности различают пять типов
конденсированных систем (табл.~\ref{tab:condensed}, рис.~\ref{fig:condensed}).

\begin{table}[ht]
\centering
\caption{Типы конденсированных систем}
\label{tab:condensed}
\small
\begin{tabularx}{\textwidth}{@{}l>{\raggedright\arraybackslash}p{33mm}>{\raggedright\arraybackslash}X@{}}
\toprule
Тип & Равновесие, симметрия, порядок & Особенности \\
\midrule
Жидкости & равновесные, изотропные, неупорядоченные & текучесть $\varphi=1/\eta$ ($\eta$~— вязкость); модуль сдвига $G=0$, распространяются только продольные волны \\
Стёкла & квази\-равновесные, изотропные, неупорядоченные & механически~— твёрдые тела: $G\neq0$, есть упругость формы, распространяются и продольные, и поперечные волны \\
Аморфные тела & сильно неравновесные, изотропные, неупорядоченные & получают в экстремальных условиях; неравновесность видна по взрывчатости аморфной сурьмы, самопроизвольному растрескиванию «батавских слёзок» \\
Жидкие кристаллы & равновесные, анизотропные, частично упорядоченные & сочетают анизотропию с большой текучестью \\
Кристаллы & равновесные, анизотропные, строго упорядоченные & дальний порядок, трансляционная симметрия (глава~1) \\
\bottomrule
\end{tabularx}
\end{table}

\begin{figure}[ht]
\centering
\begin{tikzpicture}[scale=0.42]
  % жидкость/стекло — случайная плотная упаковка (детерминированный «случайный» набор)
  \begin{scope}
    \foreach \x/\y in {0.4/0.5,1.6/0.3,2.9/0.6,4.1/0.4,5.3/0.7,0.9/1.6,2.1/1.5,3.4/1.7,4.6/1.6,5.8/1.8,
       0.3/2.7,1.5/2.8,2.7/2.6,4.0/2.9,5.2/2.8,0.8/3.9,2.0/4.0,3.3/3.8,4.5/4.0,5.7/3.9,
       0.4/5.1,1.6/5.0,2.8/5.2,4.0/5.1,5.3/5.2}{\node[atom=2.6mm] at (\x,\y) {};}
    \node[slbl] at (3,-1) {жидкость, стекло};
  \end{scope}
  % жидкий кристалл — ориентированные палочки
  \begin{scope}[shift={(8,0)}]
    \foreach \x/\y/\a in {0.5/0.6/80,1.8/0.4/85,3.0/0.7/78,4.3/0.5/88,5.5/0.6/82,
       1.0/2.4/86,2.3/2.2/79,3.6/2.5/84,4.8/2.3/90,0.4/4.2/83,1.7/4.4/88,3.0/4.1/80,4.3/4.3/85,5.6/4.0/81,5.9/2.4/87}{
      \draw[line width=2.4pt, atomA!70!black, line cap=round] ($(\x,\y)+(\a:0.7)$) -- ($(\x,\y)+(\a+180:0.7)$);}
    \node[slbl] at (3,-1) {жидкий кристалл};
  \end{scope}
  % кристалл
  \begin{scope}[shift={(16,0)}]
    \foreach \x in {0,...,5}{\foreach \y in {0,...,5}{\node[atom=2.6mm] at (\x*1.1+0.2,\y*1.0+0.2) {};}}
    \node[slbl] at (3,-1) {кристалл};
  \end{scope}
\end{tikzpicture}
\caption{Схема расположения частиц: в жидкости и стекле нет дальнего
порядка (есть только ближний); в жидком кристалле упорядочена ориентация
молекул, но не их положения; в кристалле есть и ближний, и дальний порядок}
\label{fig:condensed}
\end{figure}
